\section{Trabajo relacionado}
\label{sec:relacionado}

El problema de atender elementos según su prioridad se modela con el tipo
abstracto \emph{cola de prioridad}, cuyas operaciones básicas son insertar
un elemento y extraer el de mayor prioridad~\cite{cormen,goodrich}. La
literatura distingue entre el tipo abstracto y su representación: el mismo
comportamiento puede implementarse con estructuras distintas, y es la
representación la que determina el costo de cada operación.

Una representación directa es una lista ordenada por prioridad, que permite
extraer el elemento más prioritario en tiempo constante, pero cuya inserción
cuesta $O(n)$, pues debe recorrer la lista hasta encontrar la posición
correcta. En el extremo opuesto, una lista sin orden inserta en tiempo
constante pero obliga a recorrerla completa, en $O(n)$, cada vez que se
extrae el elemento más prioritario~\cite{goodrich}. Para prioridades
arbitrarias, la representación clásica que equilibra ambos costos es el
\emph{heap} binario, con inserción y extracción en $O(\log n)$~\cite{cormen}.

Cuando la prioridad solo puede tomar un número pequeño y fijo de valores,
como ocurre en los esquemas de triage que clasifican a cada paciente en
unos pocos niveles de urgencia, existe una alternativa más simple: mantener
una cola FIFO independiente por cada nivel. Insertar consiste en encolar en
el nivel correspondiente, y extraer consiste en revisar los niveles en orden
de urgencia hasta encontrar uno no vacío. Como el número de niveles es una
constante, ambas operaciones cuestan $O(1)$, y el orden FIFO de cada cola
resuelve de forma natural el desempate por orden de llegada. Esta es la
representación que adopta este proyecto (Sección~\ref{sec:metodologia}).

Sistemas reales de mesa de ayuda y de triage hospitalario resuelven un
problema equivalente: priorizar solicitudes por urgencia y permitir su
consulta puntual. Este proyecto no pretende replicar un sistema de
producción, sino aislar y resolver correctamente el núcleo algorítmico
común a este tipo de sistemas.
